% from ega2-1-fr.tex
\oldpage[II]{187}
de $B$ dont l'image canonique dans $B_t$ soit $\mathfrak{K}.B_t$ est l'image
réciproque de $\mathfrak{K}.B_t$, c'est-à-dire l'ensemble des $s\in B$ tels
que pour un entier $n>0$, on ait $t^ns\in\mathfrak{K}.B$. Il faut montrer que
cette dernière relation entraîne $s\in\mathfrak{K}.B$, ou encore que l'image
canonique de $t$ n'est pas diviseur de 0 dans
$B/\mathfrak{K}B=(A/\mathfrak{K})\otimes_A B$, ce qui résulte de
(\hyperref[II.8.10.2-fr]{8.10.2}) appliqué à
$X'=\operatorname{Spec}(A/\mathfrak{K})$.
\end{proof}

\begin{corollary}[8.10.5]
\label{II.8.10.5-fr}
Supposons vérifiées les hypothèses de (\hyperref[II.8.10.3-fr]{8.10.3})~;
soient $W=V-j(X)$, $x$ un point de $X$, $\mathcal{K}$ un Idéal
quasi-cohérent de $\mathcal{O}_X$, $z$ le point générique de la composante
irréductible de $\pi^{-1}(x)$ contenant $j(x)$
(\hyperref[II.8.10.3-fr]{8.10.3}).
\begin{enumerate}
  \item[\rm{(i)}] Soit $g$ une section de $\mathcal{O}_V$ au-dessus de $V$
  telle que $g|W$ soit une section de $\mathcal{K}_W$ au-dessus de $W$
  (notations de (\hyperref[II.8.10.4-fr]{8.10.4})). Alors $g$ est une section
  de $\mathcal{K}_V$~; si de plus $g(z)\neq0$, et si, pour tout entier $m>0$,
  on désigne par $g_m^x$ le germe au point $x$ de l'image canonique $g_m$ de
  $g$ dans
  $\Gamma(X,\pi_*(\mathcal{O}_V/\mathcal{J}^{m+1}))$, alors il existe un
  entier $m>0$ tel que l'image de $g_m^x$ dans
  \[
    (\pi_*(\mathcal{O}_V/\mathcal{J}^{m+1}))_x
    \otimes_{\mathcal{O}_x}k(x)
  \]
  soit $\neq0$.
  \item[\rm{(ii)}] Supposons en outre que les conditions de
  (\hyperref[II.8.10.1-fr]{8.10.1, (iii)}) soient remplies. Alors, s'il existe
  une section $g$ de $\mathcal{K}_V$ au-dessus de $V$ telle que $g(z)\neq0$,
  il existe un entier $n\geqslant0$ et une section $f$ de
  $\mathcal{K}.\mathcal{L}^{\otimes n}=
  \mathcal{K}\otimes\mathcal{L}^{\otimes n}\subset\mathcal{L}^{\otimes n}$
  telle que $f(x)\neq0$. Si $g$ est une section de $\mathcal{J}$, on peut
  prendre $n>0$.
\end{enumerate}
\end{corollary}

\begin{proof}
\begin{enumerate}
  \item[\rm{(i)}] Comme l'Idéal de $\mathcal{O}_W$ engendré par $g|W$ est
  contenu dans $\mathcal{K}_W$ par hypothèse, l'Idéal de $\mathcal{O}_V$
  engendré par $g$ est contenu dans $\mathcal{K}_V$ par
  (\hyperref[II.8.10.4-fr]{8.10.4}), autrement dit $g$ est une section de
  $\mathcal{K}_V$. Pour démontrer la seconde assertion de (i), on peut encore
  supposer $X$ et $V$ affines et conserver les notations de
  (\hyperref[II.8.10.1-fr]{8.10.1})~; la fibre $\pi^{-1}(x)$ est alors affine
  d'anneau $B'=B\otimes_A k(x)$, et il existe dans $B'$ un élément $t'$ non
  diviseur de 0 tel que $B'=k(x)\oplus B't'$. Comme $j(x)$ est une
  spécialisation de $z$ et que $g(z)\neq0$, on a nécessairement
  $g_{j(x)}\neq0$. Mais $\mathcal{O}_{j(x)}$ est un anneau local séparé
  (\hyperref[II.8.10.3-fr]{8.10.3}), qui se plonge donc dans son complété, et
  l'image de $g$ dans ce complété n'est par suite pas nulle. Or, ce complété
  est isomorphe à $\varprojlim_n(B'/B't'^{n+1})$
  (\hyperref[II.8.10.3-fr]{8.10.3})~; si $g'=g\otimes1\in B'$, il existe donc
  un entier $m$ tel que $g'\notin B't'^{m+1}$, ou encore l'image $g'_m$ de
  $g'$ dans $B'/B't'^{m+1}$ n'est pas nulle. Mais comme $g'_m$ n'est autre que
  l'image de $g_m^x$, notre assertion est démontrée.
  \item[\rm{(ii)}] En vertu de
  (\hyperref[II.8.10.1-fr]{8.10.1, (iii)})
  $\pi_*(\mathcal{O}_V/\mathcal{J}^{m+1})$ est isomorphe à la somme directe
  des $\mathcal{L}^{\otimes k}$ pour $0\leqslant k\leqslant m$~; désignons
  par $f_k$ la section de $\mathcal{L}^{\otimes k}$ au-dessus de $X$,
  composante de l'élément de
  $\bigoplus_{k=0}^m\Gamma(X,\mathcal{L}^{\otimes k})$ qui correspond à $g_m$
  par cet isomorphisme. Choisissant $m$ comme dans (i), il y a donc un indice
  $k$ tel que $f_k(x)\neq0$ d'après (i). Pour voir que $f_k$ est une section
  de $\mathcal{K}\mathcal{L}^{\otimes k}$, il suffit de considérer comme
  ci-dessus le cas où $X$ et $V$ sont affines, et cela résulte aussitôt alors
  de ce que $g\in\mathfrak{K}.B$ (notations de
  (\hyperref[II.8.10.4-fr]{8.10.4})). La dernière assertion résulte de ce que
  l'hypothèse $g\in\Gamma(V,\mathcal{J})$ entraîne $f_0=0$.
\end{enumerate}
\end{proof}

\begin{env}[8.10.6]
\label{II.8.10.6-fr}
\emph{Démonstration de (\hyperref[II.8.9.4-fr]{8.9.4}).} La question est
locale sur $Y$ (\hyperref[II.4.6.4-fr]{4.6.4})~; comme $\varepsilon$ est une
$Y$-section, on peut donc remplacer $C$ par un voisinage ouvert affine $U$
d'un point de $\varepsilon(Y)$ tel que $\varepsilon(Y)\cap U$ soit fermé dans
$U$. Autrement dit, on peut supposer $C$ affine, et $Y$ sous-préschéma fermé
de $C$ (donc affine) défini par un faisceau quasi-

